\chapter{System Requirements Specification}

\section{Introduction}
This chapter defines software requirements for the implemented Stage 1 system and the integration constraints for future Stage 2 enhancement.

\section{System Overview}
DeepSceneLoc Stage 1 is a supervised scene classification system that takes an outdoor image as input and predicts one among five scene labels: Coastal, Forest, Mountain, Rural, and Urban. The output includes probability scores and visual artifacts for evaluation.

The system is designed as a foundational semantic module in a two-stage architecture. Therefore, requirements include not only current classification behavior, but also compatibility obligations for Stage 2 location refinement.

\section{Functional Requirements}
\subsection{Input and Data Handling}
\begin{itemize}
\item The system shall ingest labeled image datasets organized by class folders.
\item The system shall apply deterministic train/validation/test splitting.
\item The system shall support image resizing and normalization compatible with pretrained backbones.
\end{itemize}

\subsection{Training}
\begin{itemize}
\item The system shall support transfer learning using ResNet-50.
\item The training module shall log epoch-level loss, accuracy, and learning rate.
\item The system shall store model checkpoints during training.
\item The system shall allow configurable hyperparameters through project configuration files.
\end{itemize}

\subsection{Evaluation and Reporting}
\begin{itemize}
\item The system shall compute overall and class wise metrics.
\item The system shall generate confusion matrices and training curves.
\item The system shall export evaluation outputs in machine-readable format.
\end{itemize}

\subsection{Inference}
\begin{itemize}
\item The system shall accept a single input image and return predicted class with confidence.
\item The demo interface shall present class probabilities for interpretability.
\end{itemize}

\subsection{Model Governance}
\begin{itemize}
\item The system shall preserve experiment metadata for each training run.
\item The system shall retain at least one best-checkpoint artifact per validated run.
\item The system shall store confusion matrix outputs for post-run analysis.
\end{itemize}

\section{Non-Functional Requirements}
\begin{itemize}
\item \textbf{Performance:} The system should support practical inference latency on commodity GPU/CPU.
\item \textbf{Reliability:} Training and evaluation outputs should be reproducible under fixed seeds and splits.
\item \textbf{Scalability:} Data pipeline should support extension to larger datasets and additional classes.
\item \textbf{Maintainability:} Source code should remain modular across data, model, evaluation, and utility layers.
\item \textbf{Usability:} A lightweight interface should allow non-expert users to test inference behavior.
\end{itemize}

\section{Primary Use-Case Scenario}
\subsection{Use-Case: Scene Inference for Unknown Image}
\textbf{Actor:} User or evaluator

\textbf{Preconditions:}
\begin{itemize}
\item model checkpoint is available,
\item inference module is active,
\item input image is readable.
\end{itemize}

\textbf{Main flow:}
\begin{itemize}
\item User uploads image.
\item System validates format and dimensions.
\item System applies preprocessing transformations.
\item Model generates logits and class probabilities.
\item System returns top prediction and confidence distribution.
\end{itemize}

\textbf{Alternate flow:}
\begin{itemize}
\item If image is invalid, system returns structured error and asks for a valid image.
\item If checkpoint is missing, system enters fallback mode or aborts with diagnostic message.
\end{itemize}

\textbf{Postconditions:}
\begin{itemize}
\item result is displayed,
\item inference event may be logged for analysis.
\end{itemize}

\section{Hardware and Software Requirements}
\subsection{Hardware}
\begin{itemize}
\item GPU-enabled training environment preferred.
\item Minimum 8 GB RAM and sufficient storage for dataset and checkpoints.
\end{itemize}

\subsection{Software}
\begin{itemize}
\item Python 3.8 or higher.
\item PyTorch and TorchVision for model development.
\item NumPy, scikit-learn, Matplotlib for metrics and visualization.
\item Gradio for demonstration interface.
\end{itemize}

\section{Constraints}
\begin{itemize}
\item Limited project duration in minor phase.
\item Incomplete landmark-level ground truth for exact location prediction.
\item Compute and API usage limits for Stage 2 experimentation.
\end{itemize}

\section{Assumptions}
\begin{itemize}
\item Input images belong to outdoor scenes relevant to the five target classes.
\item Training and test distributions are reasonably aligned after preprocessing.
\item Pretrained model weights are available for transfer learning.
\end{itemize}

\section{Future Stage Requirements}
For Stage 2, the following additional requirements are anticipated:
\begin{itemize}
\item Exact-location hypothesis generation (landmark/city/country level).
\item Confidence calibration between scene and location outputs.
\item Robust fallback behavior when fine-grained confidence is low.
\end{itemize}

\section{Acceptance Criteria}
The Stage 1 system is considered acceptable when the following conditions are met:
\begin{itemize}
\item reportable macro metrics are generated from held-out evaluation,
\item confusion matrix is available and interpretable,
\item inference demo provides class probabilities,
\item training artifacts and logs are reproducible,
\item code modules remain organized and maintainable.
\end{itemize}

\section{Risk Controls}
Risk controls embedded in requirements include deterministic split strategy, explicit checkpoint persistence, and multi-metric reporting to avoid misleading single-number conclusions. These controls support consistent progression toward Stage 2 without destabilizing Stage 1 baseline quality.

\section{Chapter Summary}
The SRS formalizes a reliable baseline system with clear extension points for advanced geolocation capabilities in the next phase.

